% abstract, choose between abstract and summary


Within the branch of Anti-Forensics, secure deletion has emerged as a critical component in cybersecurity,
regulatory compliance and digital Forensics. While standard deletion operations remove filesystem
references without erasing the underlying data,  secure deletion techniques aim to ensure that data
is made irrecoverable. This Thesis presents an experimental framework for evaluating the effectiveness, performance, and observability 
of secure deletion techniques in modern solid-state drives. The framework addresses the lack of standardized 
and reproducible evaluation methods by considering measurable factors such as post-deletion data 
recoverability, deletion latency, timing variability and verification limitations. The experimental study  
evaluates data recoverability after TRIM operations at different time intervals, and examines the execution 
and verification of drive-level sanitization commands. The experiments also examine how filesystem journaling, 
overwrite operations under Flash Translation Layer remapping, drive fill level, and vendor-specific controller 
behavior influence deletion outcomes. The results are analyzed to identify the reliability, limitations, and security 
implications of current deletion mechanisms, and to determine how effectively they support verifiable and 
standards-aligned data sanitization. In addition to the technical evaluation, this thesis also explores the 
regulatory and legal implications of secure deletion practices, particularly in the context of data protection 
laws and individual privacy rights. By establishing a reproducible evaluation methodology, this work contributes a 
foundation for systematically assessing secure deletion techniques and bridging the gap between technical 
capabilities and regulatory compliance requirements.

The importance of secure deletion is highlighted in cases where sensitive 
information, such as personal data, financial records, or confidential business information, expected 
by users to be permanently removed, can be recovered by malicious actors or forensic investigators. This 
disconnect between user expectation and technical reality has led to a cascade of high-profile data 
breaches across public and private sectors. Further the widespread transition from HDDs to SSDs has
resulted in a paradigm shift in data storage and management, introducing new challenges for secure deletion
through out-of-place updates, wear levelling, garbage collection, and Flash Translation Layer abstraction. 
These mechanisms can preserve residual copies of deleted data even after deletion or overwriting operations 
have been performed.


\vspace{0.3cm}

\textbf{Keywords} : \textit{secure data deletion; NAND flash memory; solid-state drives; Flash Translation Layer; data remanence; data sanitization; cryptographic erasure; deletion verification; regulatory compliance.}